% 34 x 34 inch workshop poster for COLM 2026
% Paper: Role-Specialized Mixture-of-Agents with Open-Weight LLMs
%        for Clinical Prediction
%
% Compile with:
%   pdflatex DAIH_COLM_2026_Poster_34x34.tex
%
% Optional files (the poster still compiles without them):
%   figures/vt-logo-white.pdf
%   figures/daih-logo-white.pdf

\documentclass[final,xcolor=table]{beamer}

% -----------------------------------------------------------------------------
% Poster dimensions: beamerposter custom dimensions are in centimetres.
% 34 inches = 86.36 cm.
% -----------------------------------------------------------------------------
\usepackage[size=custom,width=86.36,height=86.36,scale=1.0]{beamerposter}

% Only packages used by this poster.
\usepackage[T1]{fontenc}
\usepackage{lmodern}
\usepackage{amsmath}
\usepackage{array}
\usepackage{booktabs}
\usepackage{graphicx}
\usepackage{hyperref}
\usepackage{qrcode}
\usepackage{ragged2e}
\usepackage{tikz}
\usetikzlibrary{arrows.meta,calc,positioning}

\hypersetup{colorlinks=true,urlcolor=VTMaroon}

% -----------------------------------------------------------------------------
% Colour system
% -----------------------------------------------------------------------------
\definecolor{VTMaroon}{HTML}{861F41}
\definecolor{VTOrange}{HTML}{E87722}
\definecolor{DeepNavy}{HTML}{16324F}
\definecolor{AnalystBlue}{HTML}{2B6F91}
\definecolor{ProtectTeal}{HTML}{2A7F7B}
\definecolor{IntegratorGold}{HTML}{F2A900}
\definecolor{SoftGray}{HTML}{F3F5F7}
\definecolor{MidGray}{HTML}{5A6570}
\definecolor{PaleOrange}{HTML}{FFF2E7}
\definecolor{PaleBlue}{HTML}{EAF3F8}
\definecolor{PaleTeal}{HTML}{E8F5F3}

\setbeamercolor{background canvas}{bg=white}
\setbeamercolor{normal text}{fg=DeepNavy,bg=white}
\setbeamercolor{headline}{fg=white,bg=VTMaroon}
\setbeamercolor{block title}{fg=white,bg=VTMaroon}
\setbeamercolor{block body}{fg=DeepNavy,bg=SoftGray}
\setbeamercolor{takeaway}{fg=white,bg=VTOrange}
\setbeamercolor{resultbox}{fg=DeepNavy,bg=PaleOrange}
\setbeamercolor{notebox}{fg=DeepNavy,bg=PaleBlue}
\setbeamercolor{boundarybox}{fg=DeepNavy,bg=PaleTeal}

\setbeamertemplate{navigation symbols}{}
\setbeamertemplate{blocks}[rounded][shadow=false]
\setbeamertemplate{itemize item}{\color{VTOrange}$\blacktriangleright$}
\setbeamertemplate{itemize subitem}{\color{VTOrange}$\bullet$}
\setbeamersize{text margin left=1.15cm,text margin right=1.15cm}

% Physical font sizes for the printed poster.
\setbeamerfont{title}{size=\fontsize{55}{61}\selectfont,series=\bfseries}
\setbeamerfont{author}{size=\fontsize{25}{29}\selectfont,series=\bfseries}
\setbeamerfont{institute}{size=\fontsize{20}{24}\selectfont}
\setbeamerfont{block title}{size=\fontsize{30}{35}\selectfont,series=\bfseries}
\setbeamerfont{block body}{size=\fontsize{24}{29}\selectfont}
\setbeamerfont{normal text}{size=\fontsize{24}{29}\selectfont}

\newcommand{\tightitems}{%
  \setlength{\itemsep}{0.35em}%
  \setlength{\parskip}{0pt}%
  \setlength{\parsep}{0pt}%
}
\newcommand{\smallcapslabel}[1]{\textcolor{VTMaroon}{\textbf{#1}}}
\newcommand{\metric}[2]{\textbf{#1}\;\textcolor{VTMaroon}{\textbf{#2}}}
\newcommand{\optionalgraphic}[2]{%
  \IfFileExists{#1}{\includegraphics[#2]{#1}}{}%
}

% -----------------------------------------------------------------------------
% Metadata
% -----------------------------------------------------------------------------
\mode<presentation>{\usetheme{Lankton}}

\addtobeamertemplate{headline}{}
{
  \begin{tikzpicture}[remember picture,overlay]
    \node[anchor=north east,inner sep=3cm]
    at ([xshift=0.8cm,yshift=1.4cm]current page.north east)
    {
      \includegraphics[height=6cm]
      {Newposter-PromptLink/figures/EMNLP2025_main-logo_sand.png}
    };
  \end{tikzpicture}
}

% Only these two fields are changed
\title{Role-Specialized Mixture-of-Agents with Open-Weight LLMs
\\[10pt]
for Clinical Prediction}

\author{Jun Hou, Yi Fang, Xuan Wang}

% Unchanged from the original template
\institute{{junh, yif, xuanw}@vt.edu}


\begin{document}
\begin{frame}[t]

% =============================================================================
% One-sentence story
% =============================================================================
\begin{beamercolorbox}[wd=\textwidth,rounded=true,sep=0.65em]{takeaway}
  \centering
  {\fontsize{28}{34}\selectfont\bfseries
  Large open-weight analysts plus a small open-weight integrator match Claude 3.5
  few-shot on MIMIC-IV mortality Macro-F1 (58.8 vs. 57.4) while increasing
  sensitivity (61.1 vs. 15.8), without task-specific training.}
\end{beamercolorbox}

\vspace{0.55cm}

% =============================================================================
% Top row: motivation + central method figure
% =============================================================================
\setlength{\columnsep}{0.75cm}
\begin{columns}[t,totalwidth=\textwidth]

  \begin{column}{0.30\textwidth}
    \begin{block}{Motivation and Questions}
      \justifying
      Clinical prediction with closed-model APIs raises privacy and compliance
      concerns. Local fine-tuning avoids external APIs but requires labeled data,
      training resources, and task-specific adaptation.

      \vspace{0.45em}
      Existing medical multi-agent systems are commonly evaluated as a single
      block, obscuring which role controls the final prediction.

      \vspace{0.55em}
      \begin{beamercolorbox}[wd=\linewidth,rounded=true,sep=0.50em]{notebox}
        \textbf{Q1.} Which role determines the operating point at which patients
        are flagged?\\[0.35em]
        \textbf{Q2.} Can open-weight agents approach closed-model prompting
        without training?
      \end{beamercolorbox}

      \vspace{0.50em}
      \smallcapslabel{Our design:} two label-blind contrastive analysts and one
      final integrator, all served locally.
    \end{block}
  \end{column}

  \begin{column}{0.685\textwidth}
    \begin{block}{Role-Specialized Open-Weight Pipeline}
      \centering
        \begin{figure}[t]
        \centering
        \includegraphics[width=0.32\linewidth]{Newposter-PromptLink/figures/Paper/main.pdf}
        \caption{Pipeline overview. \textbf{(A) Multi-source databases.} A knowledge database is built from MedCorp and UMLS, and a patient database is built from the EHRs. \textbf{(B) Role-specialized retrieval and integration.} The label-blind risk analyst and protective analyst read the target record and augment their reasoning through retrieval from the knowledge database and the patient database; the integrator runs its own retrieval and integrates the two readings into a final prediction.}
        \label{fig:pipeline}
        \end{figure}

      \vspace{0.20em}
      \begin{columns}[T,totalwidth=0.97\linewidth]
        \begin{column}{0.32\linewidth}
          \textcolor{AnalystBlue}{\textbf{Contrastive reasoning}}\\
          Each analyst receives one outcome-specific exemplar but does not see
          its label or make an early prediction.
        \end{column}
        \begin{column}{0.32\linewidth}
          \textcolor{VTOrange}{\textbf{Controlled retrieval}}\\
          Retrieval resources and exemplar pools remain fixed when topology and
          model placement change.
        \end{column}
        \begin{column}{0.32\linewidth}
          \textcolor{VTMaroon}{\textbf{Training-free deployment}}\\
          Both model families run locally; no task-specific fine-tuning or closed
          API is required.
        \end{column}
      \end{columns}
    \end{block}
  \end{column}

\end{columns}

\vspace{0.45cm}

% =============================================================================
% Bottom region: three columns
% =============================================================================
\begin{columns}[t,totalwidth=\textwidth]

% -----------------------------------------------------------------------------
% LEFT COLUMN
% -----------------------------------------------------------------------------
\begin{column}{0.295\textwidth}

  \begin{block}{Evaluation Design}
    \begin{itemize}\tightitems
      \item \textbf{Data:} MIMIC-III and MIMIC-IV, approximately 1,000 test
        patients per dataset-task cell.
      \item \textbf{Outcomes:} in-hospital mortality and 15-day readmission.
      \item \textbf{Models:} GPT-OSS-120B (large) and Qwen2.5-7B (small), both
        served locally.
      \item \textbf{Controlled factors:} topology, medical-knowledge retrieval,
        patient exemplars, and role assignment.
      \item \textbf{Metrics:} Macro-F1, sensitivity, specificity, AUROC, and AUPRC.
    \end{itemize}

    \vspace{0.35em}
    \centering
    \begin{tabular}{@{}lcc@{}}
      \toprule
      \textbf{Task on MIMIC-IV} & \textbf{$n$} & \textbf{Positive rate} \\
      \midrule
      Mortality & 996 & 19.2\% \\
      Readmission & 996 & 46.5\% \\
      \bottomrule
    \end{tabular}
  \end{block}

  \vspace{0.35cm}

  \begin{block}{What Carries the Gain?}
    \smallcapslabel{Retrieved context alone is insufficient.}
    For the large model on MIMIC-IV mortality:

    \vspace{0.45em}
    \centering
    \begin{tabular}{@{}>{\raggedright\arraybackslash}p{0.65\linewidth}r@{}}
      \toprule
      \textbf{Configuration} & \textbf{F1} \\
      \midrule
      CoT, no retrieval & 58.3 \\
      Single + RAG & 37.7 \\
      Single + RAG + exemplars & 52.8 \\
      Multi-agent + exemplars & 51.3 \\
      Multi-agent + RAG + exemplars & \textbf{54.4} \\
      \bottomrule
    \end{tabular}

    \vspace{0.55em}
    \justifying
    RAG can hurt a single agent. Patient exemplars become useful when the two
    neighbours are divided across contrastive analyst roles. Removing MedRAG
    from the all-large pipeline changes F1 by only 3.1 points.

    \vspace{0.45em}
    \begin{beamercolorbox}[wd=\linewidth,rounded=true,sep=0.48em]{notebox}
      \centering\textbf{The organization of evidence across roles carries more
      of the gain than medical-knowledge retrieval itself.}
    \end{beamercolorbox}
  \end{block}

\end{column}

% -----------------------------------------------------------------------------
% MIDDLE COLUMN
% -----------------------------------------------------------------------------
\begin{column}{0.405\textwidth}

  \begin{block}{Main Result: High-Recall Mortality Prediction}
    \centering
    \textbf{MIMIC-IV mortality, full test set, threshold $0.5$}

    \vspace{0.40em}
    \rowcolors{2}{white}{SoftGray!70}
    \begin{tabular}{@{}>{\raggedright\arraybackslash}p{0.47\linewidth}ccc@{}}
      \toprule
      \textbf{Training-free method} & \textbf{Accuracy} & \textbf{Macro-F1} &
      \textbf{Sensitivity} \\
      \midrule
      Claude 3.5 few-shot & 84.5 & 57.4 & 15.8 \\
      GPT-OSS single + RAG + exemplars & 71.6 & 52.8 & 22.1 \\
      All-large MoA & 78.7 & 54.4 & 14.7 \\
      \rowcolor{IntegratorGold!24}
      \textbf{Large analysts + small integrator} & \textbf{66.4} &
      \textbf{58.8} & \textbf{61.1} \\
      \bottomrule
    \end{tabular}

    \vspace{0.70em}
    \begin{columns}[T,totalwidth=0.98\linewidth]
      \begin{column}{0.48\linewidth}
        \begin{beamercolorbox}[wd=\linewidth,rounded=true,sep=0.65em]{resultbox}
          \centering
          {\fontsize{31}{35}\selectfont\bfseries +1.4 F1}\\[0.15em]
          relative to Claude 3.5 few-shot
        \end{beamercolorbox}
      \end{column}
      \begin{column}{0.48\linewidth}
        \begin{beamercolorbox}[wd=\linewidth,rounded=true,sep=0.65em]{resultbox}
          \centering
          {\fontsize{31}{35}\selectfont\bfseries +45.3 sensitivity}\\[0.15em]
          relative to Claude 3.5 few-shot
        \end{beamercolorbox}
      \end{column}
    \end{columns}

    \vspace{0.65em}
    \justifying
    The mixed assignment also exceeds the sensitivity of GRU, AdaCare, and
    GraphCare by 3.3--18.2 points, although those supervised baselines retain
    substantially higher F1.
  \end{block}

  \vspace{0.35cm}

  \begin{block}{The Benefit Is an Operating-Point Shift}
    \begin{columns}[T,totalwidth=\linewidth]
      \begin{column}{0.48\linewidth}
        \begin{beamercolorbox}[wd=\linewidth,rounded=true,sep=0.55em]{notebox}
          \centering
          \textcolor{AnalystBlue}{\textbf{Large integrator}}\\[0.25em]
          Conservative threshold\\
          High specificity\\
          More missed moderate-risk deaths
        \end{beamercolorbox}
      \end{column}
      \begin{column}{0.48\linewidth}
        \begin{beamercolorbox}[wd=\linewidth,rounded=true,sep=0.55em]{resultbox}
          \centering
          \textcolor{VTMaroon}{\textbf{Small integrator}}\\[0.25em]
          Lower threshold\\
          High sensitivity\\
          More true positives and false positives
        \end{beamercolorbox}
      \end{column}
    \end{columns}

    \vspace{0.60em}
    \begin{itemize}\tightitems
      \item Changing only the integrator raises the predicted-positive rate by
        \textbf{30.2 points} on MIMIC-IV mortality.
      \item Both model families associate more recorded conditions with higher
        risk, but Qwen crosses the positive threshold earlier.
      \item \textbf{56\%} of decedents missed by the GPT-OSS integrator are
        flagged by the Qwen integrator.
    \end{itemize}

    \vspace{0.35em}
    \centering
    \begin{tabular}{@{}lccc@{}}
      \toprule
      \textbf{Full-test model} & \textbf{Single large} & \textbf{All large} &
      \textbf{Mixed} \\
      \midrule
      AUROC & 0.763 & 0.730 & 0.706 \\
      \bottomrule
    \end{tabular}

    \vspace{0.35em}
    \justifying
    \textbf{Scope:} the mixed system improves case finding at the default cutoff;
    it does not uniformly improve ranking discrimination.
  \end{block}

\end{column}

% -----------------------------------------------------------------------------
% RIGHT COLUMN
% -----------------------------------------------------------------------------
\begin{column}{0.275\textwidth}

  \begin{block}{Role Isolation: The Integrator Decides}
    \centering
    \textbf{MIMIC-IV mortality, fixed $n=100$ subset}

    \vspace{0.40em}
    \begin{tabular}{@{}lccc@{}}
      \toprule
      \textbf{A/R/I} & \textbf{F1} & \textbf{Sens.} & \textbf{AUROC} \\
      \midrule
      oss/oss/oss & 38.5 & 13.0 & .543 \\
      oss/qwen/oss & 35.0 & 7.4 & .533 \\
      qwen/oss/oss & 36.3 & 11.1 & .517 \\
      \rowcolor{IntegratorGold!24}
      \textbf{oss/oss/qwen} & \textbf{57.4} & \textbf{64.8} & \textbf{.603} \\
      \bottomrule
    \end{tabular}

    \vspace{0.45em}
    \justifying
    Replacing only the analyst or retrieval model does not produce the gain.
    Replacing the integrator changes the final operating point.

    \vspace{0.35em}
    {\fontsize{18}{22}\selectfont
    A/R/I = analyst/retrieval/integrator; oss = GPT-OSS-120B;
    qwen = Qwen2.5-7B.}
  \end{block}

  \vspace{0.35cm}

  \begin{block}{Why Readmission Is an Honest Negative}
    \begin{beamercolorbox}[wd=\linewidth,rounded=true,sep=0.50em]{boundarybox}
      \centering
      \textbf{Available-record feature AUROC}\\[0.25em]
      Mortality: \textbf{0.77}\qquad Readmission: \textbf{0.59}
    \end{beamercolorbox}

    \vspace{0.45em}
    The lower threshold transfers to readmission, but the pre-discharge record
    contains much weaker signal for this post-discharge outcome:

    \vspace{0.35em}
    \begin{itemize}\tightitems
      \item Sensitivity: $43.5 \rightarrow 94.5$
      \item Macro-F1: $47.2 \rightarrow 34.7$
      \item Specificity falls to approximately $4\%$
    \end{itemize}

    \vspace{0.35em}
    \begin{beamercolorbox}[wd=\linewidth,rounded=true,sep=0.48em]{boundarybox}
      \centering\textbf{A sensitive integrator cannot recover outcome signal
      that is absent from the available record.}
    \end{beamercolorbox}
  \end{block}

  \vspace{0.35cm}

  \begin{block}{Takeaways}
    \begin{enumerate}\tightitems
      \item \textbf{Role placement matters:} the integrator controls the
        high-recall operating point.
      \item \textbf{Topology matters:} contrastive role separation makes
        retrieved context more usable.
      \item \textbf{Task signal matters:} mortality benefits, while readmission
        exposes the boundary of the approach.
    \end{enumerate}

    \vspace{0.45em}
    \begin{columns}[c,totalwidth=\linewidth]
      \begin{column}{0.47\linewidth}
        \centering
        \qrcode[height=3.80cm]{https://arxiv.org/abs/2608.22176}\\[-0.05em]
        {\fontsize{18}{22}\selectfont\textbf{Paper}}
      \end{column}
      \begin{column}{0.47\linewidth}
        \centering
        \qrcode[height=3.80cm]{https://github.com/JuneHou/moa-clinical-rag}\\[-0.05em]
        {\fontsize{18}{22}\selectfont\textbf{Code}}
      \end{column}
    \end{columns}
  \end{block}

\end{column}

\end{columns}

\vspace{0.45cm}

% =============================================================================
% Scope strip: uses the remaining square-poster width and keeps claims calibrated
% =============================================================================
\begin{beamercolorbox}[wd=\textwidth,rounded=true,sep=0.58em]{notebox}
  \textcolor{VTMaroon}{\textbf{Scope and limitations}}\par\vspace{0.25em}
  \begin{columns}[T,totalwidth=0.98\linewidth]
    \begin{column}{0.235\linewidth}
      \textbf{Evidence concentration.} The mixed assignment shows clear
      discrimination primarily on MIMIC-IV mortality.
    \end{column}
    \begin{column}{0.235\linewidth}
      \textbf{Generalizability.} Evaluation is limited to MIMIC and two model
      families, without external or prospective validation.
    \end{column}
    \begin{column}{0.235\linewidth}
      \textbf{Retrieval design.} Each patient receives one exemplar from each
      outcome class; top-$k$ retrieval was not evaluated.
    \end{column}
    \begin{column}{0.235\linewidth}
      \textbf{Clinical review.} Intermediate reasoning was not reviewed by
      clinicians, and additional clinical outcomes remain untested.
    \end{column}
  \end{columns}
\end{beamercolorbox}

\vfill

% =============================================================================
% Compact footer
% =============================================================================
{\fontsize{15}{18}\selectfont
\begin{beamercolorbox}[wd=\textwidth,sep=0.42em]{notebox}
  \textbf{Selected references:}
  Jiang et al., ``KARE,'' WWW 2025.
  \quad Jin et al., ``MedCPT,'' \emph{Bioinformatics} 2023.
  \quad Wang et al., ``Mixture-of-Agents,'' ICLR 2025.\\[0.18em]
  \textbf{Acknowledgments:}
  NSF 2442253 and 2607580; NIH 1R21AG091260-01; USDA NIFA; Commonwealth Cyber
  Initiative; Nvidia; Cisco; Amazon--Virginia Tech; and Delta through
  ACCESS/NAIRR allocation NAIRR240202.
\end{beamercolorbox}
}

\end{frame}
\end{document}
